\documentclass[10pt]{article}
\usepackage[T1]{fontenc}
\usepackage[utf8]{inputenc}
\usepackage[french]{babel}
\usepackage{amsmath,amssymb,mathrsfs,mathpazo}
\usepackage{geometry}
\geometry{a4paper,margin=2cm}
\usepackage{graphicx,array,tabularx,tabularray,colortbl,multirow,multicol,booktabs,enumitem,pifont,ulem,fancybox,diagbox,cancel}
\usepackage{pstricks,pstricks-add,pst-plot,pst-tree,pst-node,pst-text,pst-func,pst-math,pst-bspline,pst-3dplot,pst-eucl,pst-solides3d}
\let\pscalculate\undefined
\usepackage{pst-bezier}
\usepackage{tikz,pgfplots,tkz-tab}
\usetikzlibrary{arrows,arrows.meta,patterns,calc,positioning,trees,decorations.markings,shapes.geometric,shapes.misc,shapes.arrows,intersections,decorations.pathmorphing,decorations.pathreplacing,backgrounds,shadows,babel}
\pgfplotsset{compat=1.18}
\usepgfplotslibrary{fillbetween}
\usepackage[np]{numprint}
\usepackage{listings,makecell,pas-tableur,twoopt}
\let\barre\undefined
\usepackage[pstricks]{tablvar}
\definecolor{darkgreen}{RGB}{0,100,0}
\providecommand{\red}{\color{red}}
\providecommand{\blue}{\color{blue}}
\providecommand{\green}{\color{green}}
\providecommand{\black}{\color{black}}
\providecommand{\white}{\color{white}}
\providecommand{\gray}{\color{gray}}
\providecommand{\cyan}{\color{cyan}}
\providecommand{\magenta}{\color{magenta}}
\providecommand{\yellow}{\color{yellow}}
\providecommand{\violet}{\color{violet}}
\providecommand{\orange}{\color{orange}}
\providecommand{\euro}{\texteuro}
\providecommand{\EUR}{\texteuro}
\providecommand{\ieme}{\textsuperscript{e}}
\providecommand{\degres}{\ensuremath{^\circ}}
\providecommand{\degree}{\ensuremath{^\circ}}
\providecommand{\vv}[1]{\overrightarrow{#1}}
\providecommand{\Coord}[1]{#1}
\providecommand{\decosix}{}
\providecommand{\starredbullet}{\ensuremath{\bullet}}
\pagestyle{empty}
\setlength{\parindent}{0pt}
\usepackage{adjustbox}\DeclareUnicodeCharacter{2236}{:}\DeclareUnicodeCharacter{2190}{\ensuremath{\leftarrow}}
\providecommand{\R}{\mathbb{R}}
\providecommand{\N}{\mathbb{N}}
\providecommand{\Z}{\mathbb{Z}}
\providecommand{\Q}{\mathbb{Q}}
\providecommand{\D}{\mathbb{D}}
\providecommand{\C}{\mathbb{C}}
\providecommand{\e}{{\mathrm{e}}}
\providecommand{\vect}[1]{\overrightarrow{#1}}
\providecommand{\vectt}[1]{\overrightarrow{\mathrm{#1}}}
\providecommand{\vvt}[1]{\overrightarrow{\mathrm{#1}}}
\providecommand{\np}[2][]{{#2~#1}}
\providecommand{\vv}[1]{{\overrightarrow{#1}}}
\renewcommand{\d}{{\mathrm{d}}}
\providecommand{\degres}{{^\circ}}
\providecommand{\degre}{{^\circ}}
\providecommand{\blue}{\color{blue}}
\providecommand{\red}{\color{red}}
\providecommand{\enskip}{\quad}
\providecommand{\llbracket}{\lbrack\!\lbrack}
\providecommand{\rrbracket}{\rbrack\!\rbrack}
\providecommand{\euro}{{\text{€}}}
\providecommand{\no}{{\text{n°}}}
\providecommand{\interval}[2]{{[#1;#2]}}
\providecommand{\intervFO}[2]{{[#1;#2[}}
\providecommand{\dotfill}{\ldots}

\providecommand{\ds}{\displaystyle}
\providecommand{\textding}[1]{#1}
\providecommand{\Oij}{\(\mathrm{O};\vec{\imath},\vec{\jmath}\)}
\providecommand{\Oijk}{\(\mathrm{O};\vec{\imath},\vec{\jmath},\vec{k}\)}
\providecommand{\Ouv}{\(\mathrm{O};\vec{u},\vec{v}\)}
\providecommand{\intervFF}[2]{\ensuremath{[#1;#2]}}
\renewcommand{\euro}{\ifmmode\text{\texteuro}\else\texteuro\fi}
\DeclareUnicodeCharacter{00B0}{\ifmmode^{\circ}\else\textdegree\fi}
\DeclareUnicodeCharacter{20AC}{\ifmmode\text{\texteuro}\else\texteuro\fi}
% Énoncés : Éducation nationale. Transcription et corrigés : APMEP, Denis Vergès et contributeurs. Découpage et classement : CoopMaths, Sébastien Lozano, Éric Elter et contributeurs. Adaptation de la présentation : Maths Explorer.
% CC BY-SA : https://coopmaths.fr/www/about/
% https://coopmaths.fr/alea/static/bac/2023/tex/bac_2023_03_sujet1_ameriquenord_4.tex
\begin{document}
\noindent\textit{APMEP / CoopMaths. CC BY-SA. Présentation adaptée par Maths Explorer.}\par\medskip
\medskip

On considère la suite $\left(u_{n}\right)$ définie par $u_{0}=5$ et pour tout entier naturel $n$,

\[	u_{n+1}=\frac{1}{2}\left(u_{n}+\frac{11}{u_{n}}\right)\]

On admet que la suite $\left(u_{n}\right)$ est bien définie.

\medskip

\textbf{Partie A - Étude de la suite \boldmath$ \left(u_{n}\right)$}\unboldmath

\medskip

\begin{enumerate}
\item Donner $u_{1}$ et $u_{2}$ sous forme de fractions irréductibles.
\item On considère la fonction $f$ définie sur l'intervalle $]0 ; +\infty [$ par :

\[f(x)=\frac{1}{2}\left(x+\frac{11}{x}\right)\]

Démontrer que la fonction $f$ est croissante sur l'intervalle $\left[\sqrt{11}~;~+\infty\right[$.
\item Démontrer par récurrence que pour tout entier naturel $n$, on a $: u_{n} \geqslant u_{n+1} \geqslant \sqrt{11}$.
\item En déduire que la suite $\left(u_{n}\right)$ converge vers une limite réelle. On note $a$ cette limite.
\item Après avoir déterminé et résolu une équation dont $a$ est solution, préciser la valeur exacte de $a$.
\end{enumerate}

\medskip

\textbf{Partie B - Application géométrique}

\medskip

Pour tout entier naturel $n$, on considère un rectangle $R_{n}$ d'aire 11 dont la largeur est notée $\ell_{n}$ et longueur $L_{n}$

La suite $\left(L_{n}\right)$ est définie par $L_{0}=5$ et, pour tout entier naturel $n$,

\[L_{n+1}=\frac{L_{n}+\ell_{n}}{2} \]

\begin{enumerate}
\item 
	\begin{enumerate}
		\item Expliquer pourquoi $\ell_{0}=2,2$.
		\item Établir que pour tout entier naturel $n$,

\[\ell_{n}=\frac{11}{L_{n}}.\]
	\end{enumerate}
\item Vérifier que la suite $\left(L_{n}\right)$ correspond à la suite $\left(u_{n}\right)$ de la \textbf{partie A}.
\item Montrer que pour tout entier naturel $n$, on a $\ell_{n} \leqslant \sqrt{11} \leqslant L_{n}$.
\item On admet que les suites $\left(L_{n}\right)$ et $\left(\ell_{n}\right)$ convergent toutes les deux vers $\sqrt{11}$. Interpréter géométriquement ce résultat dans le contexte de la \textbf{partie B}.
\item Voici un script, écrit en langage Python, relatif aux suites étudiées dans cette partie :

\begin{center}
\begin{tabular}{|>{\bfseries \ttfamily}l >{\ttfamily}l|} \hline
1 & def heron(n): \\
2 & \quad $L$=5 \\
3 & \quad $\ell$=2.2 \\
4 & \quad for i in range(n): \\
5 & \qquad $L$ = ($L$+$\ell$) / 2 \\
6 & \qquad $\ell$ = 11 / $L$ \\
7 & \quad return round($\ell$,6), round($L$,6)\\
\hline
\end{tabular}
\end{center}

On rappelle que la fonction Python \texttt{round(x,k)} renvoie une version arrondie du nombre \texttt{x} avec \texttt{k} décimales.

	\begin{enumerate}
		\item Si l'utilisateur tape \texttt{heron(3)} dans une console d'exécution Python, qu'obtient-il comme valeurs de sortie pour $\ell$ et $L$ ?
		\item Donner une interprétation de ces deux valeurs.
	\end{enumerate}
\end{enumerate}

\end{document}
